% ============================================================================
% Section 4: EEPROM Driver (Task 4)
% Handout requires: calculation of A, concise write/read transaction sequence,
% status mechanism used to detect write completion, write/read result,
% CS/SCK oscilloscope capture.
% ============================================================================
\section{EEPROM Driver}
\label{sec:eeprom}

\paragraph{Group address.}
\begin{equation}
  A = (n_1 + 3n_2) \bmod 256 = (\nOne + 3 \times \nTwo) \bmod 256
    = \ASum \bmod 256 = \ADec = \text{\AHex}.
\end{equation}

\paragraph{Address format.}
The fitted EEPROM takes a two-byte address, most significant byte first, and
holds 8192 bytes, not the single address byte of the CAT25040 named in the
handout~\cite{readme}.
\writeup{How you confirmed the address format on the bench (e.g. what happened to a write/read with the wrong number of address bytes, and the WEL bit afterwards).}

\paragraph{Transaction sequences.}
Every transaction is framed by CS low \ldots\ CS high and built from single-byte
\reg{spi\_transfer()} calls. Opcodes and status bits are from the EEPROM
datasheet~\cite{eeprom}.

\begin{table}[htbp]
  \centering
  \caption{EEPROM transactions used by the driver.}
  \label{tab:transactions}
  \small
  \begin{tabularx}{\linewidth}{@{}l l X r@{}}
    \toprule
    \textbf{Operation} & \textbf{Opcode} & \textbf{Bytes sent (MOSI) $\rightarrow$ received (MISO)} & \textbf{Clocks} \\
    \midrule
    Write enable (WREN) & \fillin{0x..} & WREN & \fillin{8} \\
    Write byte (WRITE)  & \fillin{0x..} & WRITE, A[15:8], A[7:0], data & \fillin{32} \\
    Read status (RDSR)  & \fillin{0x..} & RDSR, dummy $\rightarrow$ status & \fillin{16} \\
    Read byte (READ)    & \fillin{0x..} & READ, A[15:8], A[7:0], dummy $\rightarrow$ data & \fillin{32} \\
    \bottomrule
  \end{tabularx}
\end{table}

\writeup{Concise write sequence and read sequence as a numbered list (read status, WREN as its own transaction, WRITE, poll status until ready, read back, compare). State when the write-enable latch is set and when the internal write actually starts, with the datasheet reference.}

\paragraph{Write-completion detection.}
\writeup{Which status-register bit indicates a write in progress, its polarity, the mask used, how the polling loop is bounded (EEPROM\_WRITE\_TIMEOUT\_MS, timeout counter), and why a fixed delay is not used.}

\paragraph{Result.}
\begin{table}[htbp]
  \centering
  \caption{Write / read / verify result (debugger variables, \reg{RUN\_TASK}~$=4$).}
  \label{tab:task4-result}
  \begin{tabular}{@{}lll@{}}
    \toprule
    \textbf{Variable} & \textbf{Value} & \textbf{Meaning} \\
    \midrule
    \reg{eeprom\_test\_addr}     & \AHex & address $A$ \\
    \reg{eeprom\_test\_byte}     & \BHex & byte $B$ written \\
    \reg{eeprom\_status\_before} & \fillin{0x..} & \fillin{WEL / busy state} \\
    \reg{eeprom\_status\_after}  & \fillin{0x..} & \fillin{WEL / busy state} \\
    \reg{eeprom\_write\_wait\_ms} & \fillin{..} & write-cycle time (datasheet max \fillin{..~ms}) \\
    \reg{eeprom\_read\_value}    & \fillin{0x..} & byte read back \\
    \reg{eeprom\_verify\_ok}     & \fillin{1} & read value equals $B$ \\
    LEDs PB7\ldots PB0           & \fillin{bbbbbbbb} & read value on LED7\ldots LED0 \\
    \bottomrule
  \end{tabular}
\end{table}

\paragraph{Persistence test.}
\writeup{Write B with PA0, reset / power-cycle the board, and report what the read-only path (runs at reset, never writes) showed on the LEDs and in eeprom\_read\_value.}

\scopefig{task4_cs_sck.png}{Real EEPROM transaction: CH1 = CS (EEPROM pin~1 /
  PB12), CH2 = SCK (PB13). \fillin{Name the transaction shown}:
  \fillin{n} clock pulses occur while CS is low, as expected
  (\cref{tab:transactions}).}{fig:task4}
